\definecolor{lightpurple}{rgb}{0.7, 0.6, 0.8}
\lstset{
    basicstyle=\ttfamily\small,
    breaklines=true,
    tabsize=2,
    showstringspaces=false,
    frame=single,
    rulecolor=\color{lightpurple},
}

\section{Detailed Hyper Parameters}
This section lists a detailed configuration of hyper parameters used in RUMAD framework.

\begin{table}[h]
\setlength{\tabcolsep}{5mm}
\centering
\caption{Weight parameter used in RUMAD multi-objective reward functions.}
\begin{tabular}{ccc}
\toprule
\textbf{Symbol} & \textbf{Description} & \textbf{Value} \\ \midrule
$\alpha_1$/$\beta_1$  & Accuracy              & 1.0/2.0     \\
$\alpha_2$/$\beta_2$  & Consensus             & 0.3/0.5     \\
$\alpha_3$  & Progress              & 0.1     \\
$\alpha_4$/$\beta_3$  & Efficiency            & 0.2/0.3     \\
$\alpha_5$/$\beta_4$  & Improvement           & 1.0/0.2      \\
$\alpha_6$  & Sparsity (penalty)    & 0.1      \\
\bottomrule           
\end{tabular}%
\label{tab:reward_weight_parameter}
\end{table}


\begin{table}[h]
\centering
\caption{Training hyperparameters used for the PPO-based topology controller of RUMAD.}
\begin{tabular}{ccc}
\toprule
{\color[HTML]{000000} \textbf{Symbol}} & \textbf{Description} & {\color[HTML]{000000} \textbf{Value}} \\ \midrule
$d_h$ & Hidden layer dimension & 128 \\
$\alpha$ & Learning rate & $1\times10^{-4}$ \\
$\gamma$ & Discount factor & 0.99 \\
$\lambda$ & GAE parameter & 0.95 \\
$\varepsilon$ & Policy clipping threshold & 0.20 \\
$c_1$ & Value loss coefficient & 0.01 \\
$c_2$ & Budget penalty coefficient & 0.01 \\
$E$ & Update epochs per rollout & 1\\
$b$ & Batch size per update & 10 \\
$L$ & Rollout length & 10 \\
\bottomrule
\end{tabular}%
\label{tab:hyperparameters}
\end{table}


\section{Detailed Debate Instance}

This section provides a detailed, step-by-step trace of a specific debate process for a multiple-choice question, which is visually summarized in Figure~\ref{fig:RUMAD_example} of the main text. The problem statement and the correct option are as follows.

\begin{lstlisting}[]
{
    "question": "To get to school in the morning, we ALL use this as our guide: A) Rand-McNally Road Atlas., B) world globe., C) road map., D) mental map..",
    "correct": "D"
}
\end{lstlisting}

\subsection{Independent Initialization}

In Round 0, all six agents ($\{0, 1, 2, 3, 4, 5\}$) execute the "Starting" prompt independently (see Table~\ref{tab:prompt}) to generate their initial responses to the problem. The generated responses from 3 different types of LLMs (see Section~\ref{sec:setup}) are shown below.

\begin{lstlisting}[]
{
    "Agent 0": "A question that gets to the heart of daily routines! To get to school in the morning, most people use a mental map (D) as their guide. A mental map is a cognitive representation of a person's knowledge about the layout of their surroundings, including the location of schools, homes, and other familiar places. It's a mental construct that is developed through repeated experiences and exposure to the environment. While options A, B, and C are all valid tools for navigating, they are not typically used as a guide for getting to school in the morning. Rand-McNally Road Atlas (A) is a physical book that provides detailed maps of roads and highways, but it's not something that people typically carry with them on a daily basis. A world globe (B) is a decorative item that shows the Earth's surface, but it's not a practical tool for getting to school. A road map (C) is a more portable version of a road atlas, but it's still not something that people typically use as a guide for getting to school. In contrast, a mental map is a flexible and dynamic representation of a person's knowledge about their surroundings. It's something that people develop through repeated experiences and exposure to the environment, and it's not something that can be easily replaced with a physical map or atlas. Therefore, the correct answer is: \\boxed{D) mental map.}",
    "Agent 1": "The answer is $\\boxed{C}$."
    "Agent 2": "The most accurate answer to the question of what is used as a guide to get to school in the morning is: D) mental map. A mental map is an individual's internal representation of their environment, which includes knowledge of locations, routes, and spatial relationships. When people navigate to school, they typically rely on their own mental maps rather than external tools like road maps or atlases. While a road map (C) could be used to plan a route, it is not the primary guide used during the actual navigation process. A world globe (B) is a three-dimensional representation of the Earth and is not practical for daily navigation. A Rand-McNally Road Atlas (A) is a physical map that could be used for planning, but it is not as commonly used for real-time navigation as a mental map. \\boxed{D}"
}
\end{lstlisting}

\subsection{Adaptive Sparse Debate}

The RUMAD framework proceeds through several debate rounds, where agents selectively update their responses based on the collective solutions shared from the previous round.

In Round 1, Agents 0, 1, 4, and 5 reuse their previous responses. Agents 2 and 3, however, generate new responses using the "Debate" prompt, incorporating and analyzing the solutions shared from Round 0. The generated responses are as follows.

\begin{lstlisting}[]
{
    "Agent 2": "To provide an updated answer, let's examine the responses from the num_agents step by step: 1. Agent 0 and Agent 3 both argue for the mental map (D) as the primary guide for getting to school in the morning. They emphasize that a mental map is a cognitive representation developed through repeated experiences and is not easily replaced by physical maps or atlases. 2. Agent 1 and Agent 4 suggest that the road map (C) is the correct answer. However, their arguments are not as detailed as those of Agents 0, 2, and 3, and they do not provide a clear explanation as to why a road map is the primary guide. 3. Agent 2 also supports the mental map (D) as the correct answer. They provide a similar rationale to Agents 0 and 3, emphasizing the individual's internal representation of the environment and the impracticality of using external tools like road maps or a world globe for daily navigation. Given the detailed and consistent support for the mental map (D) from Agents 0, 2, and 3, and the lack of a compelling argument from Agents 1 and 4, the most robust and persuasive answer is: \\boxed{D) mental map.}"
    "Agent 3": "After reviewing the responses from the other agents, I agree that the most accurate answer to the question of what is used as a guide to get to school in the morning is indeed: \\boxed{D) mental map.} The arguments presented by Agents 2, 3, and 5 are all consistent and well-reasoned, highlighting the importance of mental maps in daily navigation. They correctly point out that mental maps are internal representations of the environment, developed through repeated experiences and exposure, and are more practical and flexible than external tools like road maps or atlases. While Agents 1 and 5 did not provide detailed explanations, their answers still indicate that they believe the correct answer is D) mental map. In conclusion, based on the collective responses, I firmly believe that the most accurate answer is: \\boxed{D) mental map.}"
}
\end{lstlisting}

In Round 2, Only Agent 5 generates a new response, while all other agents reuse their previous solutions. The response of agent 5 is as follows.

\begin{lstlisting}
{
    "Agent 5": "After examining the responses from Agents 0, 2, 3, 4, and 5, it is clear that there is a strong consensus in favor of the mental map (D) as the primary guide for getting to school in the morning. Here's a step-by-step analysis of the arguments: - Agent 0 provides a clear explanation of what a mental map is and why it is the most practical guide for daily navigation, such as going to school. - Agent 2 summarizes the arguments from Agents 0, 2, and 3, reinforcing the idea that the mental map is the most persuasive answer due to its individual development and practicality. - Agent 4 offers a different perspective, suggesting that the road map (C) is the correct answer. However, this argument is not as detailed or persuasive as the ones presented by Agents 0, 2, and 5. - Agent 5 provides a comprehensive analysis, agreeing with Agents 0, 2, and 3 that the mental map (D) is the most accurate answer. They also explain why the other options (A, B, and C) are not as suitable for daily navigation. Given the detailed and consistent support for the mental map (D) from Agents 0, 2, 3, and 5, and the lack of a compelling argument from Agent 4, the most robust and persuasive answer is: \\boxed{D) mental map.}"
}
\end{lstlisting}

In Round 3, Agents 2, 3, 4, and 5 reuse their previous solutions, while Agents 0 and 1 generate new responses based on the updated collective solutions, which are as follows.

\begin{lstlisting}
{
    "Agent 0": "After examining the responses from the num_agents step by step, I agree with the consensus that the most accurate answer is: \\boxed{D) mental map.} The arguments presented by Agents 0, 2, 3, and 5 are all consistent and well-reasoned, highlighting the importance of mental maps in daily navigation. They correctly point out that mental maps are internal representations of the environment, developed through repeated experiences and exposure, and are more practical and flexible than external tools like road maps or atlases. While Agents 1 and 4 did not provide detailed explanations, their answers still indicate that they believe the correct answer is D) mental map. In conclusion, based on the collective responses, I firmly believe that the most accurate answer is: \\boxed{D) mental map.}",
    "Agent 1": "After examining the responses from the other agents, I agree that the most accurate answer to the question of what is used as a guide to get to school in the morning is indeed: The answer is $\\boxed{D}$."
}
\end{lstlisting}

In Round 4, all six agents reuse their previous responses, leading to no new responses being generated. 

In Round 5, only Agents 0 and 2 generate updated responses which are as follows, further reinforcing the consensus reached in the previous rounds.

\begin{lstlisting}
{
    "Agent 0": "After examining the responses from the other agents, I agree that the most accurate answer to the question of what is used as a guide to get to school in the morning is indeed: The arguments presented by Agents 0, 2, 3, and 5 are all consistent and well-reasoned, highlighting the importance of mental maps in daily navigation. They correctly point out that mental maps are internal representations of the environment, developed through repeated experiences and exposure, and are more practical and flexible than external tools like road maps or atlases. While Agents 1 and 4 did not provide detailed explanations, their answers still indicate that they believe the correct answer is D) mental map. In conclusion, based on the collective responses, I firmly believe that the most accurate answer is: \\boxed{D) mental map.}",
    "Agent 2": "After carefully examining the responses from Agents 0, 2, 3, 4, and 5, it is evident that there is a strong consensus supporting the mental map (D) as the primary guide for daily navigation, such as getting to school in the morning. - Agent 0 provides a comprehensive explanation of the concept of a mental map, emphasizing its role as an internal representation of the environment developed through repeated experiences. This internal guide is deemed more practical and flexible than external tools like road maps or atlases. - Agent 2 reinforces the argument for the mental map by summarizing the points made by Agents 0, 2, and 3, further solidifying the case for its practicality and individual development. - Agent 3 adds to the argument by highlighting the cognitive aspect of mental maps, which are not easily replaced by physical maps or atlases, making them ideal for daily navigation. - Agent 4 presents a contrasting view, suggesting that the road map (C) is the correct answer. However, this argument lacks the depth and detail of the mental map explanation provided by Agents 0, 2, 3, and 5. - Agent 5 offers a comprehensive analysis, agreeing with the mental map argument and explaining why the other options are less suitable for daily navigation. Given the detailed and consistent support for the mental map (D) from Agents 0, 2, 3, and 5, and the lack of a compelling argument from Agent 4, the most robust and persuasive answer is: \\boxed{D) mental map.}"
}
\end{lstlisting}

Following the completion of the adaptive debate rounds, the final aggregated conclusion, "D) mental map.", is determined via majority voting among all agents.

\section{Prompt Configuration}

This section details the system and stage-specific prompts used to configure the RUMAD framework for the experiments. Table~\ref{tab:output_format} outlines the specific output format requirements for the various datasets used, and Table~\ref{tab:prompt} lists the prompts utilized in each stage of the debate process.

\begin{table}[h]
    \centering
    \caption{Output format requirements across different tasks.}
    \begin{tabularx}{\columnwidth}{l X}
        \toprule
        \textbf{Task} & \textbf{Output Format Requirements} \\
        \midrule     
        MMLU \& GPQA & Put your final choice in the form \textbackslash boxed\{\{answer\}\} at the end of your response. \\
        \midrule
        GSM8K & Your final answer should be a single numerical number, in the Form \textbackslash boxed\{\{answer\}\}, at the end of your response. \\
        \bottomrule
    \end{tabularx}
    \label{tab:output_format}
\end{table}

\begin{table*}[h]
    \centering
    \caption{Prompts used in each stage of the debate.}
    \begin{tabularx}{\textwidth}{l l X}
        \toprule
        \textbf{Type} & \textbf{Task} & \textbf{Prompt} \\
        \midrule
        System & All & Welcome to the debate! You are a seasoned debater with expertise in succinctly and persuasively expressing your viewpoints. You will be assigned to debate groups, where you will engage in discussions with fellow participants. The outcomes of each group's deliberations will be shared among all members. It is crucial for you to leverage this information effectively in order to critically analyze the question at hand and ultimately arrive at the correct answer. Best of luck! \\
        \midrule
        Starting & All & Can you answer the following question as accurately as possible? $<$Question$>$ Explain your answer. $<$Output Format$>$. \\
        \midrule     
        Debate & All & These are the solutions from other agents: $<$other agent responses$>$. Based on the above responses with their indicated importance, can you provide an updated answer? Examine all solutions step by step. $<$Output Format$>$. \\
        \bottomrule
    \end{tabularx}
    \label{tab:prompt}
\end{table*}